% MAPA artigo_3:5.4
\begin{proposition}[forma ancorada e comprimento estrutural]\label{map:artigo_3:5.4}
\label{artigo_3:prop:anchored-relu-code}
Se $g\in\Aclass^\square_{N,H}([0,1])$, sua forma reduzida possui
$M\le N-1$ hinges e pode ser escrita, tomando o centro $x_*$ da folha mais à
esquerda, como
\begin{equation}
 g(x)=g(x_*)+\sigma_0(x-x_*)
      +\sum_{j=1}^{M}(\sigma_j-\sigma_{j-1})\rho(x-\tau_j).
 \label{artigo_3:eq:anchored-relu-form}
\end{equation}
Os valores $g(x_*)$ e $\sigma_i$ têm altura no máximo $H$, os saltos têm
altura no máximo $2H^2$, e os pontos diádicos são transmitidos por seus pares
estruturais. Consequentemente, existe uma codificação prefix-free da forma
ancorada com
\begin{equation}
 L_{\rm ReLU}^{\rm anc}(g)=O(N\log(2H)).
 \label{artigo_3:eq:anchored-network-bits}
\end{equation}

Na direção contrária, se os coeficientes de
\eqref{artigo_3:eq:reduced-relu-reverse} têm altura no máximo $W$ e as quebras têm
profundidade no máximo $D$ (e as quebras são diádicas em $(0,1)$), então
\begin{equation}
 H_{\rm can}(R)\le 2^{D+1}\,(M+2)\,W^{M+3},
 \qquad
 \log_2 H_{\rm can}(R)
 \le 5\bigl[(M+1)\log_2(2W)+D+\log_2(M+1)\bigr].
 \label{artigo_3:eq:reverse-height-bit-bound}
\end{equation}
\end{proposition}

\begin{proof}
Na ida, cada salto é a diferença de duas inclinações de altura no máximo
$H$, logo sua altura não excede $2H^2$. A altura atômica limita os inteiros
estruturais da célula, e todos os campos usam $O(\log(2H))$ bits. Somar os
$O(N)$ campos prova \eqref{artigo_3:eq:anchored-network-bits}, com constante explícita:
o código formalizado usa no máximo $40\,N\lfloor\log_2(2H)\rfloor$ bits. O ancoramento evita
expandir desnecessariamente $g(x_*)-\sigma_0x_*$ em um intercepto com
denominador diádico explícito.

Na volta, toda folha $J_{k,m}$ de $\mathcal T(S)$ tem $k\le D$ e
$|m|+1\le 2^k\le 2^D$. Cada inclinação local $\sigma_i$ é soma de no máximo $M+1$
racionais de altura $\le W$: um denominador comum é $\le W^{M+1}$ e
$|\sigma_i|\le(M+1)W$, logo $\operatorname{ht}\sigma_i\le (M+1)W^{M+2}$. O centro
$x_{k,m}=(2m+1)/2^{k+1}$ e as quebras têm denominador dividindo $2^{D+1}$, e
$0\le \max\{0,x_{k,m}-\tau_j\}\le1$. Assim, $R(x_{k,m})=c+qx_{k,m}+\sum_j
a_j\max\{0,x_{k,m}-\tau_j\}$ admite o denominador comum
$\operatorname{den}(c)\operatorname{den}(q)\,2^{D+1}\prod_j\operatorname{den}(a_j)\le
2^{D+1}W^{M+2}$ (a parte diádica é compartilhada por todas as parcelas) e módulo
$\le (M+2)W$; portanto $\operatorname{ht}R(x_{k,m})\le 2^{D+1}(M+2)W^{M+3}$. Os
demais termos de \eqref{artigo_3:eq:canonical-network-height} ($2$, $k+1\le D+1$,
$|m|+1\le 2^D$) são menores, o que dá $H_{\rm can}(R)\le 2^{D+1}(M+2)W^{M+3}$.
Tomando $\log_2$, e usando $W\ge1$, $\log_2(M+2)\le 1+\log_2(M+1)$ e
$(M+1)\log_2(2W)\ge1$, obtém-se
\[
 \log_2 H_{\rm can}\le D+1+\bigl(1+\log_2(M+1)\bigr)+(M+3)\log_2 W
 \le 5\bigl[(M+1)\log_2(2W)+D+\log_2(M+1)\bigr].
\]
\end{proof}

\begin{observacaoSN}\label{artigo_3:rem:reverse-height-correction}
O fator é $M+1$, e não $M$, no primeiro termo: há $M+1$ inclinações locais e
cada uma é soma de até $M+1$ racionais. A versão com $M\log(2W)$ é falsa já para
$M=0$: em $R(x)=x$ ($M=0$, $D=0$, $W=1$) tem-se $H_{\rm can}=2$, de modo que
$\log_2 H_{\rm can}=1$ e o lado direito é $0$; e em $R(x)=Wx$ tem-se
$H_{\rm can}\ge W$ (a altura da inclinação), de modo que nenhum termo constante
repara a cota. A volta corrigida está formalizada em
\texttt{Pendencias/AlturaInversa.lean}; a ida (forma ancorada e código prefix-free)
está formalizada em \texttt{Pendencias/FormaAncorada.lean} e
\texttt{Pendencias/CodigoAncorado.lean}.
\end{observacaoSN}

